\chapter{The order book}\label{ch:orders}

Orders are entered once, in the order book. Each planning period takes the open orders due in it; orders are
marked delivered when the period that delivered them is closed. You never type orders into a period by hand,
except for the rare period-only order (Section~\ref{sec:period-orders}).

\section{Page overview}

\shot{orders}{The order book with the expected load and the list of open orders.}{orders}

\begin{callouts}
  \item[1] \textbf{Expected load, next 8 weeks}: the open orders' demand per stage and week (below).
  \item[2] \textbf{Search} by order number, product or customer.
  \item[3] \textbf{Status filter}: Open (default), Delivered, Cancelled, or all; the \ui{\texttimes} clears it.
  \item[4] \textbf{Template}: downloads a CSV with the expected columns, in the interface language.
  \item[5] \textbf{Import from Excel}: imports an .xlsx/.csv file or pasted rows (Section~\ref{sec:import}).
  \item[6] \textbf{Add order}: enter a single order.
  \item[7] \textbf{Sortable columns}: click a header to sort, again to reverse.
\end{callouts}

\section{Expected load}

The panel converts the open orders' demand into machine-hours per stage, per week (using the fastest machine
and excluding setups), and compares it with the working calendar. A percentage above 85 is tight. Hover a
cell to see ``need / available machine-hours''. Above the table, two warnings may appear: how many open orders
are already overdue, and how many orders due within seven days are in no open period. Both are also shown
on the Today page.

\section{The order list}

Each row shows the order number, the product (with its colour dot), quantity, due date and time, priority,
customer and status. Overdue open orders are marked \ui{overdue}. The list is paged; the result count is
shown on the right.

\section{Adding, editing and deleting an order}

Click \ui{Add order}, or click an order number to edit it:

\shotf[0.5\linewidth]{order_dialog}{Editing an order.}{order_dialog}

\begin{center}
\begin{tabularx}{\linewidth}{>{\bfseries}p{3cm}X}
\toprule
\rowcolor{white}Field & Meaning \\
\midrule
Order code & Unique identifier (for example \code{SO-261124-039} or a customer PO). Fixed after creation. \\
Product & One of the factory's products. \\
Quantity & Units ordered (a whole number). \\
Due & Date and time the customer expects the goods (plant local time). \\
Priority & \ui{Low} (1), \ui{Normal} (2), \ui{High} (3), \ui{Urgent} (4). Lateness is weighted by priority, so an hour late on an urgent order costs four times an hour on a low one. \\
Pickup & \ui{At due time}: the customer collects at the due time, so goods finished early wait in the warehouse (and count against the stock cap). \ui{When ready}: the order ships as soon as it is finished. \\
Customer, Note & Free text, shown in periods and exports. \\
Status & \ui{Open}, \ui{Delivered} or \ui{Cancelled}. Closing a period sets delivered orders automatically; you rarely change it by hand. \\
\bottomrule
\end{tabularx}
\end{center}

Changes are saved to the book immediately. Open periods that use the order notice the change and offer to
take it in (Section~\ref{sec:sync}).

\begin{warning}
\ui{Delete} removes the order from the book completely, and periods using it are prompted to update. To keep
a record of an order the customer withdrew, set its status to \ui{Cancelled} instead.
\end{warning}

\section{Importing orders from Excel, CSV or a paste}\label{sec:import}

New orders usually arrive as a spreadsheet. The import reads it, shows a preview of every row with the
problems it found, and writes nothing until you confirm.

\shotf[0.82\linewidth]{import_start}{Step 1: choose a file or paste rows copied from Excel.}{import_start}

\begin{steps}[Import orders]
\begin{enumerate}
  \item Click \ui{Import from Excel} (also in the \ui{New} menu and on the Today page).
  \item Either \ui{Choose an .xlsx or .csv file} (the first sheet is read; the first row holds the headers),
        or copy rows from Excel, paste them in the text box and click \ui{Preview the pasted rows}.
  \item Read the preview (Figure~\ref{fig:import_preview}): the badges summarise what will happen, the
        \ui{Columns found} line shows how each header was understood, and each row has a result.
  \item Choose the mode, then click \ui{Write \textit{n} changes}.
\end{enumerate}
\end{steps}

\shotf[0.85\linewidth]{import_preview}{Step 2: the preview. Three new orders, one update of an existing order, one row with an unknown product.}{import_preview}

\subsection{How the file is read}

\begin{itemize}
  \item \textbf{Headers} may be English or Vietnamese, with or without accents, in any order: ``Order no'',
        ``PO'', ``Mã đơn'' all mean the order code; ``Qty'', ``Số lượng'' the quantity; and so on.
        Appendix~\ref{app:import} lists every accepted header. Unknown columns are ignored and named in
        the preview.
  \item \textbf{Required columns}: order code, product, quantity and due date. If one is missing, the preview
        says ``Missing required columns'' and nothing can be written.
  \item \textbf{Dates}: \code{dd/mm/yyyy} (also \code{dd-mm-yy} or \code{dd.mm.yyyy}), ISO
        \code{yyyy-mm-dd}, or real Excel dates. A time may follow the date or sit in its own column
        (\code{16:00}, \code{16h}); without a time, the factory's default due time is used.
  \item \textbf{Quantities} are whole numbers; \code{1.200} or \code{1,200} are read as twelve hundred.
  \item \textbf{Priority} may be a number 1--9 or a word: Low (Thấp), Normal (Thường), High (Cao), Urgent or
        Rush (Gấp).
  \item \textbf{Pickup} words: ``at due'', ``on due'', ``đúng hẹn'' for at due time; ``when ready'',
        ``ASAP'', ``khi xong'' for when ready.
  \item \textbf{Product} must be a product code; a close misspelling is suggested in the check notes.
\end{itemize}

\subsection{Row results}

\begin{center}
\begin{tabularx}{\linewidth}{>{\bfseries}p{2.4cm}X}
\toprule
\rowcolor{white}Result & Meaning \\
\midrule
Add & A new order code: it will be created. \\
Update & The code exists and is open; the check notes list what changes (for example ``changes: Quantity, Note''). \\
Same & The code exists with identical values; nothing to do. \\
Skip & The order is already delivered or cancelled; it is never changed by an import. \\
Error & The row cannot be used; the reason is shown in red (unknown product, bad date, missing quantity\dots). \\
\bottomrule
\end{tabularx}
\end{center}

If some rows have errors, writing is blocked until you either fix the file or tick \ui{Skip the \textit{n} rows
with errors and write the rest}.

\subsection{Import modes}

\begin{description}[style=nextline, leftmargin=1.2em, font=\bfseries\color{brand}]
  \item[Add \& update] (default) Orders in the file are added or updated; all other orders stay as they are.
        Use it for a file of new or changed orders.
  \item[Replace all open orders] Open orders missing from the file are \emph{cancelled}. Use it only when the
        file is the complete list of open orders (for example an export from the sales system). The preview
        shows how many open orders will be cancelled before you confirm.
\end{description}

\section{How orders reach planning periods}\label{sec:orders-periods}

\begin{itemize}
  \item When a period is created, it takes every \textbf{open} order due before its end --- including orders
        that are already overdue, which are scheduled as early as possible --- plus the open orders due in
        the look-ahead window after it (shown as ``after the period'').
  \item Changes made later in the book (new, changed or cancelled orders) are \emph{not} pushed into periods
        silently. The period's Prepare step shows a banner and you take them in with one click
        (Section~\ref{sec:sync}).
  \item An order becomes \ui{Delivered} when the period whose executed schedule delivers it is closed. Until
        then it stays open --- which is why a committed period should be closed before the next one is
        created (Section~\ref{sec:new-period}).
\end{itemize}
